\begin{abstract}
We characterize the value of outcome-free covariate-treatment records for estimating a conditional average treatment effect at the center of the unit covariate cube. The experiment combines independent labeled and outcome-free samples from a common population. Under a known uniform covariate design, conditional exchangeability, a public overlap level \(0<\varepsilon<1/2\) with propensity values in \([\varepsilon,1-\varepsilon]\), binary responses with arm means in \([0,1]\), and public Hölder regularity, we establish the sharp minimax expected absolute-error rate uniformly over labeled counts of at least two and every nonnegative auxiliary count. The propensity and control-mean smoothness orders lie in \((0,1]\), the contrast smoothness order lies in \([1,3]\), and the common public Hölder radius satisfies \(L>1/2\). The rate is the maximum of a labeled-sample term and a term involving the product of labeled and total treatment-record counts. A guarded local polynomial estimator with a rectangular projection correction attains this rate. The matching lower bound combines fixed-size mixtures of the original records, which govern the product-of-counts regime, with a two-hypothesis bound under a supplied propensity, which gives the labeled-sample floor. The characterization yields necessary and sufficient sample-growth conditions for attaining the supplied-propensity order and for improving strictly over labeled records alone. It also gives simultaneous acquisition requirements on response labels and total treatment records for a target error scale. The rate exponents do not depend on \(\varepsilon\), while the comparison constants may depend on it together with the other fixed public parameters.
\end{abstract}

\section{Introduction}\label{sec:introduction}

This paper establishes a sharp pointwise precision benchmark for conditional average treatment effect estimation when responses are observed on a smaller collection of records than covariates and treatment assignments. We consider \(n\ge2\), the labeled-record count, and \(m\ge0\), the independent outcome-free covariate-treatment record count. For the specified potential-outcome class in \cref{obj:def:primitive-class}, with known uniform covariate design and the public parameter restrictions stated below, the minimax expected absolute error is
\[
\underbrace{R(n,m)}_{\text{minimax pointwise absolute-error risk}}
\asymp
\max\left\{
n^{-\gamma/(2\gamma+d)},
[n(n+m)]^{-\gamma/(2\gamma+d+\gamma d/(\alpha+\beta))}
\right\},
\]
where \(d\) is the covariate dimension, and \(\alpha,\beta,\gamma\) are the Hölder smoothness orders of the propensity, control response mean, and causal contrast, respectively. The target is the conditional causal contrast at the cube's center. The comparison holds uniformly over both record counts, with constants depending on the fixed public class parameters. A fully specified decision based on the observed records attains this order, as established in \cref{obj:thm:sharp-annotation-frontier}.

The model fixes the covariate distribution to be uniform on \([0,1]^d\), with \(d\ge1\), and assumes conditional treatment exchangeability. Treatment and potential responses are binary. For a fixed public \(\varepsilon\in(0,1/2)\), the propensity lies in \([\varepsilon,1-\varepsilon]\), and both potential-response means lie in \([0,1]\) throughout the cube. The public smoothness orders satisfy \(\alpha,\beta\in(0,1]\) and \(\gamma\in[1,3]\), with \(L\), the common public Hölder radius, satisfying \(L>1/2\). These restrictions define the primitive law class in \cref{obj:def:primitive-class}; the independent, common-population record blocks follow \cref{obj:ass:sample-law}. Admissible conditional-margin witnesses satisfy the restrictions over the whole cube. Their continuous contrast, together with full-support design and the identifying assumptions, determines the point target uniquely from the observed law.

The distinction between record types has a direct acquisition interpretation. Response ascertainment through chart review can yield labeled records containing covariates, treatment, and the true response, while routine records supply covariates and treatment. Surrogate-assisted treatment-effect estimation provides a related electronic-record setting \citep[Sections 2.1, 2.5--2.6, and 4]{ChengAnanthakrishnanCai2021}. Within the present sampling model, adding a new independent labeled record changes \((n,m)\) to \((n+1,m)\). Labeling one already-counted outcome-free record changes the composition to \((n+1,m-1)\) at fixed total count; if the record index is fixed in advance, or chosen by a randomization independent of all records, the resulting dataset follows the stated sampling law at \((n+1,m-1)\). Adding a new outcome-free record changes \((n,m)\) to \((n,m+1)\). The sharp risk formula benchmarks these sample compositions under the stated independent-sampling assumptions.

The estimator combines a coarse local polynomial fit with a finer projection correction. It splits the original records into independent outcome and treatment roles, whose sizes are proportional to the labeled and total treatment-record counts. A rectangular average across these roles corrects the local estimating moments for the unknown nuisance functions. An empirical eigenvalue guard and clipping produce a finite Borel decision for every dataset. Under the stated model conditions, \cref{obj:thm:rectangular-upper} bounds its risk by effect approximation, nuisance approximation, localized response sampling, and rectangular sampling terms. Public tuning balances these terms to attain the sharp rate.

The matching lower bound retains the original fixed-size observation experiment. The construction in \cref{obj:thm:marked-component-certificate} couples propensity and response perturbations while preserving identical auxiliary-only mixture laws across hypotheses. Its information bound supplies the nuisance converse in \cref{obj:lem:nuisance-frontier-converse} over the range where the product-of-counts term governs precision. The supplied-propensity comparison in \cref{obj:thm:supplied-propensity} establishes the labeled-sample order \(n^{-\gamma/(2\gamma+d)}\) over the same primitive class and every auxiliary count. Together, these lower bounds and the attaining estimator give the uniform minimax characterization, including randomized decisions.

The characterization distinguishes strict improvement from attainment of the supplied-propensity order. Write \(S=\alpha+\beta\), the combined nuisance smoothness, and let \(S_{\mathrm{crit}}\) denote the smoothness threshold and \(q_*\) the total-record growth exponent:
\[
S_{\mathrm{crit}}=\frac{\gamma d}{2\gamma+d},
\qquad
q_*=\frac{\gamma d}{S(2\gamma+d)}.
\]
At and above the threshold, labeled records alone attain the supplied-propensity order. Below it, attainment along an auxiliary-size sequence requires total treatment-record counts to remain at least a positive constant times \(n^{q_*}\). Below the smoothness threshold, strict asymptotic improvement over labeled records alone occurs exactly when the auxiliary-to-labeled ratio diverges. \Cref{obj:prop:annotation-threshold} also translates a target absolute-error scale into simultaneous requirements on the labeled count and the product of labeled and total counts. These are risk requirements up to parameter-dependent constants, with exact equivalences for the displayed rate.

The closest supervised comparison is \citet[Section 3, Theorems 1--2 and equation (14)]{KennedyBalakrishnanRobinsWasserman2024}. At zero auxiliary records, our rate has the same numerical exponents and smoothness cutoff, as recorded in \cref{obj:lem:published-cate-benchmark}. Their statistical guarantees concern their bounded-density model and stated attainment conditions; the lower bound for our known-uniform class combines the original-record construction in \cref{obj:thm:marked-component-certificate,obj:lem:nuisance-frontier-converse}, which covers smoothness below \(S_{\mathrm{crit}}\) with total count at most \(n^{q_*}\), with the labeled-sample floor \(n^{-\gamma/(2\gamma+d)}\) from the constant-nuisance subfamily in \cref{obj:thm:supplied-propensity}. Semi-supervised analyses of average and quantile treatment effects characterize nuisance-dependent expansions and inference for marginal targets \citep[Theorem 2.1, Corollary 2.1, Theorem 3.1, and Corollary 3.1]{ChakraborttyDai2022}. Our contribution characterizes worst-case pointwise absolute error and its acquisition requirements for the specified conditional target. The projection construction connects local residualized regression \citep[Theorem 3]{Kennedy2023} with finite-rank second-order correction \citep[Lemma 3]{Robins2009}.



The exposition proceeds through related work, setup and assumptions, the estimator construction, sharp precision and sample requirements, and discussion and limitations. The upper-bound argument is organized in \cref{sec:upper-bound-proofs}, the original-record converse in \cref{sec:original-record-lower-bounds}, and the acquisition algebra in \cref{sec:acquisition-consequences}.
